
%%%%%%%%%%%%%%%%%%%% 前言 %%%%%%%%%%%%%%%%%%%%%

\begin{preface}

这本书的标题看上去有些故弄玄虚。如果这本书以“旁门左道”来讲线性代数，那人们为什么还要去读它呢？这本书的“旁”和“左”具体表现在哪里呢？

在回答这些问题之前，请允许我首先说明这本教材的目标读者群体。这本教材是“荣誉线性代数”课程的讲义。该课程是为数学基础较好的学生准备的{\bf 第一门}线性代数课程，面向虽还未熟悉抽象推理，但愿意超出“菜谱式”微积分课程而学习更加严谨的数学内容的学生，因而除作为线性代数入门课程外，本课程还旨在引导学生初步掌握{\bf 严谨}的数学证明与形式化定义——简言之，接触现代理论数学（抽象数学）的范式。目标读者要能借助通常在进阶教材中呈现的抽象定义和理论建构，来解释通常在入门线性代数教材中呈现的基本思想和具体实例。

本书另一特点是其编写视角并非来自代数专家，也非面向代数专家。因此我着重强调对分析学、几何学、概率论等领域重要的主题，省略了部分传统内容。例如，本书仅讨论实数域和复数域上的向量空间，完全不涉及其他域上的线性空间——因为我认为与其将时间用于介绍抽象域的概念，不如用来讲解其他学科需要掌握的经典主题。当学生后续在抽象代数课程中学习一般域的概念时，自会理解本书涉及的诸多理论结构同样适用于一般域。

此外，本书仅讨论有限维空间，且“基”均指有限基。原因在于：若不引入收敛性、范数、完备性等泛函分析基础概念，就无法对无限维空间进行深入讨论——而这显然是另一门课程（教材）的主题。因此本书不讨论无限哈默尔基：这类概念在分析和几何的大多数应用中并非必需，我认为它们更应属于抽象代数课程的范畴。

\subsection*{教学说明}

本书与通行的进阶线性代数教材存在若干差异。首先体现在基、线性无关性及生成组的定义方式上。我在书中首先将基定义为满足“任意向量都能唯一表示为线性组合”性质的向量组，随后基的两方面性质自然衍生出线性无关性（唯一性的体现）和生成性（存在性的体现）。

采用这种进路是因为，我认为基的概念远比线性无关性更重要——在大多数应用中，我们真正关心的是向量组能否构成基，而并不关注线性无关性。例如求解齐次方程组时，我们不仅需要线性无关解，更需要获取恰当数量的线性无关解（即解空间的基）。

并且，很容易向学生阐释基的重要性：基使我们能够引入坐标，从而可把研究抽象向量空间转化至研究$\mathbb{R}^n$（或$\mathbb{C}^n$）。此外，我们引入坐标是为了便于计算机运算——计算机尤其擅长矩阵处理。反观线性无关性的概念，我确实难以找到简明易懂的教学动机。

另一差异在于：我在讲解线性方程组解法之前先引入了线性变换。这种安排的缺点是直到第 \ref{chap:LinearSystems} 章才证明“仅方阵可逆”等重要结论。但是，提前定义线性变换能使行化简操作的阐述更具系统性。同时，我花费大量篇幅（两节）来阐释矩阵乘法的动机——希望充分说明这种看似奇特的乘法规则实则非常自然，且本质上是我们唯一的选择。

关于基、线性变换等的重要结论的证明（如向量空间中任意两个基包含相同数量向量这一结论）均在第 \ref{chap:LinearSystems} 章中，通过在行化简过程中数主元完成。尽管这些结论大多存在“无坐标”证明（形式上不涉及高斯消元法），但仔细分析便会发现：高斯消元与主元计数其实并未消失，只是暗含在多数证明过程中。因此，本书未采用进阶线性代数教材中常见的优雅却晦涩的“无坐标”证明，而是选用更常见于“微积分式”教材的“行化简”证明法——这样既能凸显诸证明背后的共同思路，也使数学基础有限的读者更易理解与记忆。

在第 \ref{chap:LinearSystems} 章第 \ref{sec:ChangeOfCoordinate} 节，我还提出了简单易记的基变换公式体系。

第 \ref{chap:Determinant} 章讨论行列式。我花费大量篇幅阐述行列式的来源，延后给出形式化定义。通过将行列式定义为体积计算工具，引导学生理解：若允许使用带符号体积以使行列式对各列是线性的（此时学生应已认识到线性的重要价值，且明白允许体积为负是其微不足道的代价）并假设若干自然性质，则必然推导出经典的行列式定义。需要强调的是，我并未先行设定行列式的反对称性，而是从体积的其他自然性质中推导得出。

需注意：虽然第 \ref{chap:BasicNotion} 至 \ref{chap:Determinant} 章形式上主要讨论实空间，但所有内容同样适用于复空间，甚至任意域上的向量空间。

第 \ref{chap:SpectralIntro} 章是谱理论入门，复空间$\mathbb{C}^n$在此自然登场（其形式定义虽在开篇已给出，但此前主要研究对象始终是实空间$\mathbb{R}^n$）。复特征值的出现表明：即便最初处理实矩阵（实空间上的算子），谱理论最自然的研究场域仍是复空间$\mathbb{C}^n$。本章重点在于对角化理论，并引入由特征空间构成的基的概念。

将涉及内积空间的第 \ref{chap:InnerProduct} 章置于谱理论之后，是因为我希望同步处理复数域与实数域情形，而谱理论为复空间提供了明确的动机。除却这层动机，第 \ref{chap:SpectralIntro}、\ref{chap:InnerProduct} 章内容相互独立，教师亦可选择先讲授第 \ref{chap:InnerProduct} 章。

尽管本书第 \ref{chap:SpectralAdvanced} 章介绍了若当标准形，但在单学期课程中通常没有时间涵盖该内容。我倾向于将更多时间用于第 \ref{chap:StructinInnerProdSpace}、\ref{chap:BilinearandInner} 章讨论的主题，例如正规算子与自伴算子的对角化、极分解与奇异值分解、正交矩阵的结构与取向，以及二次型理论。

我认为这些主题比若当标准形更具应用价值——尽管后者无疑具有理论美感。不过，我增设了第 \ref{chap:SpectralAdvanced} 章，因而教师能灵活调整教学内容：可选择跳过第 \ref{chap:StructinInnerProdSpace}、\ref{chap:BilinearandInner} 章部分主题，转而讲授若当分解定理。

本书（2009年版）还新增了关于对偶空间与张量的第 \ref{chap:DualandTensor} 章。我认为该部分内容（特别是张量相关章节）对一年级线性代数课程而言稍显拔高，但其中某些主题（如对偶空间中的坐标变换）可轻松纳入教学大纲。此外，该章也可作为进阶课程中张量理论的入门导引。需要说明的是，本章所有结论均适用于任意域。

在本书的呈现方式上，我力求采用非正式的表达风格：更倾向直观的几何推演而非形式化的代数运算，因此在纯粹主义者眼中本书可能不够严谨。全书通常将线性变换与其矩阵表示视为同一（只要不致引发混淆），这种简化标记法有助于初学者理解——过度强调变换及其矩阵的区别反而可能使初学者困惑。仅当这种区别至关重要时（例如分析基改变时线性变换矩阵的变化规律），我才会使用特殊标记法区分线性变换及其矩阵表示。

\end{preface}

